\lyseis{6.5}{Προβλήματα αναλογιών}
\sol{1} Ένα μπουκάλι κοστίζει $2{,}4:3=0{,}8$ €, άρα τα $8$ κοστίζουν $6{,}4$ €.
\sol{2} Ένα κιλό $6{,}5:5=1{,}3$ €, άρα $8$ κιλά κοστίζουν $10{,}4$ €. Με $20$ €: $20:1{,}3\approx15{,}4$ κιλά.
\sol{3} Ένα τεμάχιο $31{,}5:7=4{,}5$ €, άρα $12$ τεμάχια $54$ €.
\sol{4} Ο συντελεστής είναι $\frac{9}{3}=3$, άρα $y=15,\ 24,\ 36$ για $x=5,\,8,\,12$.
\sol{5} Ανάλογα: β, γ, δ. Όχι ανάλογα: α, ε ($y=x^2$) και στ (οι ηλικίες διαφέρουν πάντα κατά τα ίδια χρόνια).
\sol{6} Ανά άτομο $75$ g ρύζι και $0{,}5$ κρεμμύδι. Για $10$ άτομα: $750$ g και $5$ κρεμμύδια. Για $6$ άτομα: $450$ g και $3$ κρεμμύδια.
\sol{7} Για μία πίτα $1{,}2:6=0{,}2$ kg, άρα για $15$ πίτες $3$ kg.
\sol{8} Ταχύτητα $80$ km/h, σε $3{,}5$ ώρες $280$ km.
\sol{9} $8$ λίτρα ανά $100$ km, άρα $0{,}08$ λίτρα ανά km και $340\cdot0{,}08=27{,}2$ λίτρα.
\sol{10} Μία φωτοτυπία $0{,}07$ €, οι $120$ κοστίζουν $8{,}4$ €.
\sol{11} Ανά δέντρο $240:15=16$ kg, άρα $40$ δέντρα δίνουν $640$ kg.
\sol{12} $50\cdot1{,}1=55$ δολάρια. $132:1{,}1=120$ €.
\sol{13} Σύνολο ωρών $8$, μία ώρα $15$ €, Νίκος $45$ € και Βασίλης $75$ €.
\sol{14} $12$ ετών. Οι ηλικίες διαφέρουν πάντα κατά $3$ χρόνια, ο λόγος δεν μένει σταθερός ($\frac{2}{5}\neq\frac{12}{15}$).
\sol{15} Όχι απαραίτητα. Ισχύει μόνο αν η ταχύτητα μένει σταθερή, αλλά ο δρομέας κουράζεται και γίνεται πιο αργός.
\sol{16} Ένα κιλό $22{,}5:9=2{,}5$ €, άρα $14$ κιλά $35$ €. Με αναλογία: $\frac{9}{14}=\frac{22{,}5}{x}$, $9x=315$, $x=35$.
\sol{17} Όχι. Περισσότεροι εργάτες χρειάζονται λιγότερες ημέρες: $6$ εργάτες περίπου $3$ ημέρες. Τα ποσά αυτά δεν είναι ανάλογα, αλλά αντιστρόφως ανάλογα (επόμενη παράγραφος).
\sol{18} Ελεύθερη απάντηση.
\sol{19} Ελεύθερη απάντηση. Για διπλάσια άτομα διπλασιάζονται όλες οι ποσότητες, για μισά υποδιπλασιάζονται.
\sol{20} Ελεύθερη απάντηση.
